% Glossary terms from ch12: lines of the form \item[Term] Definition.
\item[Relative position] $\pos_{\mathrm{rel}}=\pos_B-\pos_A$, the position of agent $B$ seen from agent $A$; with the relative velocity it determines whether two agents on constant velocities collide.
\item[Relative velocity] $\vel_{\mathrm{rel}}=\vel_A-\vel_B$, the velocity of $A$ in the frame of $B$; two discs collide at time $t$ if and only if $\norm{\pos_{\mathrm{rel}}-t\vel_{\mathrm{rel}}}\le r_A+r_B$.
\item[Collision cone] $CC_{A|B}$, the set of relative velocities whose ray from the origin meets the disc of radius $r_A+r_B$ around $\pos_{\mathrm{rel}}$; a wedge of half-angle $\arcsin\bigl((r_A+r_B)/\norm{\pos_{\mathrm{rel}}}\bigr)$ bounded by the tangent rays.
\item[Velocity obstacle] $\VO_{A|B}(\vel_B)=\vel_B\oplus CC_{A|B}$, the set of velocities of $A$ that lead to a collision with $B$ if $B$ keeps its velocity; the collision cone with its apex moved to $\vel_B$.
\item[Apex (of a velocity obstacle)] The vertex of the wedge, at $\vel_B$ for a velocity obstacle, at the origin for a static obstacle and at $(\vel_A+\vel_B)/2$ for a reciprocal velocity obstacle.
\item[Tangent half-angle] $\theta=\arcsin(R/d)$ with $R=r_A+r_B$ and $d=\norm{\pos_{\mathrm{rel}}}$, the angle between the axis of the collision cone and each of its boundary rays.
\item[Time to collision] $t_c(\vel_{\mathrm{rel}})$, the first time at which the discs touch, the smaller root of $(\vel\cdot\vel)t^2-2(\pos\cdot\vel)t+(\pos\cdot\pos-R^2)=0$; infinite if the discriminant is negative or $\pos\cdot\vel\le0$.
\item[Truncated velocity obstacle] $\VO^{\ttc}_{A|B}$, the velocities that collide within the horizon $\ttc$; the wedge with its tip cut off by the disc of centre $\vel_B+\pos_{\mathrm{rel}}/\ttc$ and radius $R/\ttc$.
\item[Time horizon (velocity obstacle)] The time $\ttc$ beyond which predicted collisions are ignored by the local avoidance rule; it must exceed the decision interval and allow the next manoeuvre.
\item[Preferred velocity] $\vel^{\mathrm{pref}}$, the velocity an agent would fly if nothing were in the way, typically towards the goal or along the nominal route at the nominal speed.
\item[Feasible velocity set] The reachable velocities $\norm{\vel}\le v_{\max}$ minus the union of the truncated velocity obstacles of all obstacles; the local rule picks the feasible velocity closest to the preferred one.
\item[Non-cooperative obstacle] An object that does not react to the agent, such as an unknown drone or a bird; it is avoided with the plain velocity obstacle (full responsibility), not with a reciprocal one.
\item[Oscillation (velocity obstacles)] The step-by-step reversal of avoidance manoeuvres that occurs when two agents apply the velocity-obstacle rule simultaneously, each assuming that the other keeps its velocity; the motivation for reciprocal velocity obstacles.
